\begin{proposition}
\label{proposition-going-down-normal-integral}
Let $R \subset S$ be an inclusion of domains.
Assume $R$ is normal and $S$ integral over $R$.
Let $\mathfrak p \subset \mathfrak p' \subset R$
be primes. Let $\mathfrak q'$ be a prime of $S$
with $\mathfrak p' = R \cap \mathfrak q'$.
Then there exists a prime $\mathfrak q$
with $\mathfrak q \subset \mathfrak q'$
such that $\mathfrak p = R \cap \mathfrak q$. In other words:
the going down property holds for $R \to S$, see
Definition \ref{definition-going-up-down}.
\end{proposition}

\begin{proof}
Keep the stated primes and the original inclusion, and put
$B=S_{\mathfrak q'}$. It suffices to prove
$\mathfrak pB\cap R=\mathfrak p$, as in Lemma \ref{lemma-in-image}.
Here is the precise prime construction supplied by that equality.
The ring
$$
F=(R\setminus\mathfrak p)^{-1}(B/\mathfrak pB)
$$
is nonzero: if $1=0$ in $F$, some $r\in R\setminus\mathfrak p$
has zero image in $B/\mathfrak pB$, contradicting the equality.
A prime of $F$ pulls back through the quotient and localization
maps to a prime of $B$ containing $\mathfrak pB$ and avoiding
the image of $R\setminus\mathfrak p$. Its inverse image
$\mathfrak q$ under $S\to B$ has contraction $\mathfrak p$
in $R$. Every element of $S\setminus\mathfrak q'$ is a unit
in $B$, so none belongs to $\mathfrak q$ and
$\mathfrak q\subseteq\mathfrak q'$. This proves the desired
conclusion once the equality is established.

The inclusion $\mathfrak p\subseteq\mathfrak pB\cap R$ is
immediate. For the reverse inclusion choose an original
$z\in R$ whose image belongs to $\mathfrak pB$. Write that
image as a finite sum $\sum_{i=1}^r r_i h_i/g_i$ with
$r_i\in\mathfrak p$, $h_i\in S$ and $g_i\notin\mathfrak q'$.
Take
$$
g=\prod_{i=1}^r g_i,\qquad
y=\sum_{i=1}^r r_i h_i\prod_{k\ne i}g_k\in\mathfrak pS.
$$
Then $g\notin\mathfrak q'$ and $z=y/g$ in $B$. Since $S$
is a domain, the map $S\to B$ is injective, and the exact
identity $zg=y$ holds in $S$ itself. For the empty sum use
$g=1$ and $y=0$.

\medskip\noindent
By Lemma \ref{lemma-integral-integral-over-ideal}
there exists a monic polynomial
$P = x^m + b_{m-1} x^{m-1} + \ldots + b_0$
with $b_i\in\mathfrak p^{m-i}\subseteq\mathfrak p$
for every $0\leq i<m$, such that $P(y)=0$.
These are the full coefficient powers in the original definition
of integrality over $\mathfrak p$.

\medskip\noindent
By Lemma \ref{lemma-minimal-polynomial-normal-domain}, the
minimal polynomial of $g$ over $K=\operatorname{Frac}(R)$ has
coefficients in $R$. Write it as
$Q=x^n+\sum_{j=0}^{n-1}a_jx^j$. If every $a_j$ belonged to
$\mathfrak p$, its original equation would give
$$
g^n=-\sum_{j=0}^{n-1}a_jg^j\in\mathfrak pS
\subseteq\mathfrak p'S\subseteq\mathfrak q'.
$$
Since $n\geq1$ and $\mathfrak q'$ is prime, this would give
$g\in\mathfrak q'$, a contradiction. Therefore at least one
of the original $a_j$ is outside $\mathfrak p$.

\medskip\noindent
If $z=0$, the desired membership $z\in\mathfrak p$ already
holds. Hence assume $z\ne0$, so $z$ is a unit in the original
$K$, and retain $y=zg$. The coefficient-preserving substitution
map $K[x]\to K[x]$, $H(x)\mapsto H(zx)$, has inverse
$H(x)\mapsto H(x/z)$. Set
$$
Q'(x)=z^nQ(x/z)
=x^n+\sum_{j=0}^{n-1}z^{n-j}a_jx^j.
$$
This is monic of degree $n$ and satisfies
$Q'(y)=z^nQ(g)=0$. If $H(y)=0$ for $H\in K[x]$, then
$H(zg)=0$, so minimality of $Q$ gives $H(zx)=Q(x)D(x)$
for some $D\in K[x]$. The inverse substitution yields
$$
H(x)=Q(x/z)D(x/z)=Q'(x)\bigl(z^{-n}D(x/z)\bigr).
$$
Thus $Q'$ divides every annihilating polynomial of $y$ and is
its monic minimal polynomial; in particular $Q'$ divides the
original $P$. All coefficients of $P$ and $Q'$ belong to $R$.
Part (2) of Lemma \ref{lemma-polynomials-divide} gives
$$
z^ja_{n-j}\in\sqrt{(b_0,\ldots,b_{m-1})R}
\subseteq\mathfrak p,\qquad 1\leq j\leq n.
$$
The full powers $b_i\in\mathfrak p^{m-i}$ also yield the
stronger statement that each $z^ja_{n-j}$ is integral over
$\mathfrak p^j$, by
Lemma \ref{lemma-weighted-monic-factor-coefficients}.
Choose $a_{n-j}\notin\mathfrak p$ from the previous paragraph.
Primality first gives $z^j\in\mathfrak p$ and then
$z\in\mathfrak p$, since $j\geq1$. The required contraction
equality, and therefore the prime constructed above, follow.
\end{proof}